\documentclass[tikz,border=10pt]{standalone}
\usepackage{tikz}
\usetikzlibrary{calc}

\begin{document}
\begin{tikzpicture}
    % Define coordinates
    \coordinate (A) at (0,8);
    \coordinate (B) at (0,0);
    \coordinate (C) at (6,0);
    \coordinate (D) at (6,8);
    
    % E is on AD
    \coordinate (E) at (4.5,8);
    
    % F is on BE such that CF is perpendicular to BE
    % Line BE: y = (16/9)x
    % Line CF: y = (-9/16)(x-6)
    % Intersection: x = 486/337 approx 1.442, y = 864/337 approx 2.564
    \coordinate (F) at (1.442, 2.564);
    
    % H is on AD, vertically aligned with F
    \coordinate (H) at (1.442, 8);
    
    % G is on BC, vertically aligned with F
    \coordinate (G) at (1.442, 0);
    
    % Draw the rectangle
    \draw (A) -- (B) -- (C) -- (D) -- cycle;
    
    % Draw solid internal lines
    \draw (B) -- (E);
    \draw (E) -- (C);
    
    % Draw dashed lines
    \draw[dashed] (H) -- (G);
    \draw[dashed] (F) -- (C);
    
    % Draw right angle at G (between HG and GC)
    \draw ($(G)+(0,0.3)$) -- (G) -- ($(G)+(0.3,0)$);
    
    % Draw right angle at F (between FB and FC)
    % Draws the corner of the square opposite to F
    \draw ($(F)!3mm!(B)$) -- ($($(F)!3mm!(B)$) + ($(F)!3mm!(C)$) - (F)$) -- ($(F)!3mm!(C)$);
    
    % Labels
    \node[above left] at (A) {$A$};
    \node[above] at (H) {$H$};
    \node[above] at (E) {$E$};
    \node[above right] at (D) {$D$};
    
    \node[below left] at (B) {$B$};
    \node[below] at (G) {$G$};
    \node[below right] at (C) {$C$};
    
    \node[left] at (F) {$F$};
\end{tikzpicture}
\end{document}
